\section{Discussion, limitations, and future work}
\label{sec:lim}

\Eqsat engines mitigate the effects of
  rule ordering by non-destructively
  applying all the rules in the ruleset
  in each iteration of the algorithm.
However, in practical applications
  where saturation is unlikely,
  the engine completes only a bounded number of iterations
  before resource limits force termination.
Deciding \textit{which}
  rules to run \textit{when}
  becomes critical---splitting
  rules into batches can
  drastically alter the results.
Based on preliminary experiments,
  we find that certain strategies
  can dramatically improve the results of \eqsat tools.
One such strategy is a
  \textit{saturating scheduler}, which
  iteratively
  (1) applies saturating rules
  (\T{is\_saturating} in \autoref{fig:ruleops})
  until they reach saturation, then (2) applies
  the other rules for a single
  iteration; another is the use of
  operators like \T{compress}.

These observations call for a
  more systematic investigation
  of scheduling techniques than this paper provides, but we
  are excited to further explore
  rule scheduling in the context of \eqsat.
We have also only partially explored conditional rule inference in \enumo.
To use \enumo for inferring state-of-the-art
 rules in more complex domains like LLVM IR,
robust conditional rule inference as well as
program analyses to satisfy side conditions will be necessary.

While in this paper we used \enumo to synthesize rewrite
rules for and by equality saturation, the DSL itself is
generic and not restricted to applications in equality saturation.
\enumo lays the groundwork for future applications in
bounded model checking and sketch-guided synthesis.
For example, \enumo could be useful in axiom synthesis tools such as
  LAS~\citep{oopsla22-logic}, where the enumeration order has a large impact on the
  quality of results.

